\section{Introduction}

% What are we talking about?
Pit is a general-purpose language designed to be an alternative to C.  It is
similar to C, but with several additional features.  While the additional
features are general-purpose, they prevent buffer and integer overflow
vulnerabilities by automating memory allocation.  However if the cost to
efficiency is too great, manual memory allocation is still possible.  The
additional features also allow portability of object files after they are
compiled.

% Introduce autos
The main feature added by Pit is the ability to \emph{optionally} have memory allocation automatically
managed.  The choice as to whether to have the memory managed automatically is made for each variable.
This is accomplished by declaring one new variable ``type'', called ``auto''.  (This ``type'' is
actually a meta-type indicating that the variable is dynamically typed instead of statically typed.)
``Auto'' indicates that the variable is dynamically typed; that is, the variable will assume
the type of whatever value is assigned to it.  It also means that memory required to store its value
will be managed automatically.  (Under the hood, an auto is stored on the stack or in the data
segment as a pointer to a structure that stores the type, value, and reference count.  Most operations
on it are translated into function calls to a supporting library.)  All the other variable types are
referred to as ``primitive''.  Each of the primitive types corresponds to a variable type in C.

% Tradeoffs and choices
Clearly, using autos will often incur a performance penalty.  While there are many ways to reduce or
sometimes eliminate this overhead, it cannot be ignored.  The advantage to this approach is
that \emph{the programmer gets to choose} whether the benefit is worth the penalty, rather than being
forced to write the entire module in a low-maintenance high-overhead language or in a high-maintenance
low-overhead language.  In some parts of a module it is worth the penalty, and in some parts it is not.
It is a well-established fact that nearly all of the execution time of a program is spent in a fraction
of the code.  Reasonable estimates are usually something between 80\% and 90\% of the execution time is
spent in 10\% to 20\% of the code, which we will call the \emph{time-critical sections} of code.
This is why it is so important to let the programmer choose whether to use autos or primitive
variables:  any code that does not use primitive variables is completely protected from buffer overrun
and integer overflow vulnerabilities.  The code that has the largest impact on the execution time of
the program can still use primitive variables and manual memory management.  Allowing the programmer to
spend less time on 80\% of the code will allow him more time to meticulously audit the time-critical
sections where he chose to use primitive variables for bug, security-related or otherwise.

% Other options?
Many other languages have demonstrated how automatic memory management is beneficial to security;
however none of these are a viable replacement for C.  None allow the programmer to mix checked and
unchecked code, and none allow access to inline assembly code, which is a requirement for any
replacement for C.  In Pit, the same function may declare some variables that are managed
automatically and some that are not.

\subsection{Portability}

A long-standing problem with C is after the source code is compiled, it is no longer portable.  It is
locked into a particular version of a particular operating system on a particular CPU.  Because there
are so many combinations, this task is usually left to the end-users, contributing to the difficulty of
installation.  This is especially difficult for people who are not software developers.  This causes
much difficulty; often this by itself is enough to convince a vendor to pick one or two platforms
to release their software for, and the best hope for the rest is to use awkword compatibility layers, such as FreeBSD's
compatibility layer for Linux software.  Many programs are not available for FreeBSD at all, except
by using an executable compiled for Linux under this compatiblity layer.

Application-level software written in Pit can be portable \emph{after} it is translated into object
format (but before it is linked with the system libraries).  A program can be translated into object
format once, and the object files can be distributed
to all operating systems for a particular CPU.  During installation, the program is linked and
copied into the appropriate place on the filesystem.  Clearly, here we are talking about portability
between operating systems, not portability between CPUs.  This reduces the number of compilations
necessary to one for each CPU.


\subsubsection{detail -- move to body}

XXX: Fix this paragraph.
C's preprocessor is mostly to blame for this.  A library can declare macros that get resolved when
the calling module is compiled.  A new version of this library may change the contents of these macros.
A module written in Pit maintains the separation between the code in each module even after compilation...  mention stateless functions
Blame also rests with C's clumsyness at representing integers with a dynamic size.


Sometimes object-file portability has a cost to efficiency.  Sometimes an auto variable is required
in the API to provide abstraction for the types of the values that are passed between functions.
On the other hand, it is reasonable to rely on some values having exactly the same type on all operating systems.

For example, the number of elements in an array is a number that is absolutely guaranteed to fit in
a variable that is the size of a pointer.  Also, there is no benefit to using a variable that is
smaller than a pointer (in fact, it is usually slower unless the compiler rounds the size up), at
least when defining a module's API.  Therefore, when representing anything stored in memory, a
pointer-sized integer is the only reasonable choice, and we need no abstraction.  It is reasonable
to say that if an integer of larger size is required, the output bytecode must be different, which
means it will be necessary to recompile the module anyway.

On the other hand, consider an integer that represents an offset into a file.  A file's maximum
size is dependent on the operating system.  In the not-so-distant past, file sizes were represented
with a 32-bit variable, limiting them to 4GB.  At the time, it was wasteful to use larger values,
as an ``industrial size'' array of hard drives occupying many large machine rooms could store only a
fraction of 4GB.  Today, we use a 64-bit value for representing offsets into files, but it is
unreasonable to rely on that being large enough forever, particularly with the proliferation of
distributed storage systems such as RAID in the consumer market.  Clearly, locking a module's API into a
64-bit size for a file offset is short-sighted at best.  Standard C libraries address this by using
\emph{off\_t} to represent file offsets.  This allows the size to be increased without modifying the
source code to a program.  However, this is still unsatisfactory, as it requires the program to be
recompiled when a 128-bit integer is required.  In particular, during the transition period, the
compilation generates different assembler code depending on whether the target operating system has
switched to a 128-bit \emph{off\_t} or is still using a 64-bit \emph{off\_t}.  However, simply compiling
two versions does not work either, because other data types are likely to be in transition as well, each
of which doubles the number of compilationed versions required.  A better solution is to abstract the type of a
file offset out of the compiled code entirely.  An auto can represent the integer without
requiring a particular size; the amount of memory allocated will always be sufficient to represent the
number.




On a 32-bit architecture, a 32-bit
value will always be sufficient to represent the number of elements in an array.  If the program is
moved to a 64-bit architecture, the bytecode of the program will change, and therefore it must be
recompiled.  (We ignore compatibility modes because a 32-bit program running on a 64-bit machine
under a compatibility mode cannot have arrays larger than 4 GB.  Thus, there still cannot be more
than $2^{32}$ elements.)  Therefore, we consider 64-bit extensions to a 32-bit architecture to be a
distinct ar


Complex APIs may require auto variables in order to communicate
without being locked into an interface that cannot be extended.  For example, consider the C functions for
making network sockets.  Making a network connection requires different parameters depending on whether it is
using IPv4 or IPv6


\subsection{Comparison to Other Languages}


\begin{tabular}{ | l || c | c | c | }
\hline
\bf{Feature} & \bf{Pit} & \bf{C} & \bf{Perl} \\
\hline

\underline{Basic features} \\
\hline
Manual memory allocation & Y & Y & N \\
\hline
Automatic memory allocation and garbage collection & Y & N & Y \\
\hline
Structs & Y & Y & C \\
\hline
Inline assembly & Y & Y & N \\
\hline

\underline{High-level features} \\
\hline
Abstract data types & Y & N & Y \\
\hline
Automated compilation & Y & N & Y \\
\hline
Vector math & Y & N & N \\
\hline
\end{tabular}

Yes, No, or Clumsy

\subsubsection{Memory allocation}

The programmer has a choice.  If he declares an auto, the memory for anything stored in it is allocated
automatically, and reference-counting garbage collection is automatically performed.  If he declares a
primitive type (integer, float, pointer, string, struct, etc), allocation works the same way as it works
in C and no garbage collection is performed for that variable.  This emphasizes the choice given to the
programmer; he is free to choose either efficiency or security, portability, and maintainability.

\subsubsection{Structs}

Consider, for example, a stream of bytes received over a network socket containing non-string data, such as
integers or floats.  In C and Pit, it is easy to extract the elements from this by accessing the data with
the \emph{struct} construct.  In Perl, this is possible but clumsy; to do so, the data must be passed to the
\emph{unpack} function with a format string describing what the structure of the data is.

\subsubsection{Inline assembly}

Inline assembly code, while undesirable, is often unavoidable.  It is often necessary for low-level
interaction with hardware or for manually optimizing code.  Pit's interface to assembly code is nearly
identical to C's interface.

\subsubsection{Automated compilation}

When compiling a C program, the programmer must somehow describe each command required.  Each module of the
program must be compiled explicitly.  Then the modules must be linked with the necessary libraries explicitly
to produce the final executable file.  Additionally, for each module and library, a separate interface
description (``header file'' in C-specific terminology) must be maintained.  While tools such as \emph{make}
ease this process, it is still a complex process that the programmer must design carefully.

When compiling a Perl program, compilation starts with the main module.  When another module or library is
required, the compiler is told the name of the library so it can find it in the filesystem.  Not only does
this ease the compilation process, but it makes it easier to move the process to other systems where the
filesystem's hierarchy is different because the \emph{compiler} rather than files shipped with the source
code decides where to find dependencies.

As with Perl, each symbol in a Pit program belongs to a ``namespace''.  Namespaces may be nested within a
module, and each module implicitly defines a namespace defined by the filename of the module and the
subdirectory (if any) under which it resides in the main library directory.  Obviously
namespaces provide a separation between symbol names, preventing two or more modules from accidentally
defining conflicting symbol names, but this scheme also allows the linker to automatically find dependencies.
For example, if the main module imports a library called ``helper'', the compiler finds the interface file for
this library called ``helper.pi'' and reads the exported symbol prototypes from it.  All of the symbol
names in that library are prefixed with the word ``helper'' and a period.  Because of the prefix, when
the module is linked, the linker automatically looks for a file called ``helper.o'', ``helper.a'', etc
in the current directory and in the library search path.

While other steps may be necessary to compile or install the program, the basic compilation process is
automated.  Also, interface files are generated automatically instead of maintained manually.  These are
clear advantages over C.

\subsubsection{Abstract data types}

XXX: more work

\subsubsection{Vector math}

XXX: more work



\subsection{XXX: Junk to do}

\begin{itemize}
\item
	Target applications:
	\begin{itemize}
	\item
		Operating systems (kernel, libraries, CLI, GUI, etc).  No run-time
		interpreter is required.  Efficiency costs are optional.
		Easy interaction with in-line and external assembler code.  Easy
		interaction with existing C code.
	\item
		Threaded applications.  All features are (or will be) thread-safe.
	\item
		Network applications.  Additional features make string handling
		easier.
	\item
		Real-time systems.  Libraries and extra features are designed to
		require predictable and consistent usage of CPU and memory.
	\item
		Application software (web-based and local).  Autos improve data
		manageability in ways similar to Perl and related languages via
		easy allocation, garbage collection, hashes, managed arrays.
	\end{itemize}
\item
	Implementation
\end{itemize}
